\clearpage\thispagestyle{empty}
\pdfbookmark[0]{Quellen zur Besprechung von Bessels Schrift über den Kometen von 1807}{bessel1811-sources}
\begin{center}\textbf{Quellen zur Besprechung von Bessels Schrift über den Kometen von 1807}\end{center}
\noindent Diese KI-übersetzte Fassung in modernem Deutsch folgt Carl Friedrich Gauß, \textit{Werke}, Band~VI (Göttingen, 1874), S.~530--532. Die Besprechung erschien zuerst ohne Verfassernamen unter Königsberg in den \textit{Göttingischen gelehrten Anzeigen}, Stück 26, 16. Februar 1811, S.~250--256.
\par\smallskip\noindent In der Werkausgabe lautet der Titel von Friedrich Wilhelm Bessels Schrift \textit{Untersuchungen über die scheinbare und wahre Bahn des im Jahre 1807 erschienenen grossen Cometen}. Auf dem Titelblatt des Buchs steht \textit{Kometen}. Als Verlagsangabe nennt es Königsberg, Friedrich Nicolovius, 1810. Das digitalisierte Exemplar stammt aus der ETH-Bibliothek Zürich, Rar 33291.
\par\smallskip\noindent Die Werkausgabe druckt \textit{Monatliche Crrespondenz}, der Erstdruck \textit{Monatliche Correspondenz}. Im Erstdruck stehen die Formeln zum Gesichtsfeld untereinander, und beim Beitrag jedes einzelnen Planeten steht zusätzlich das Wort \textit{störenden}. Die Übertragung folgt der Werkausgabe.
\par\smallskip\noindent \href{https://gdz.sub.uni-goettingen.de/id/PPN236007467}{Text der Werkausgabe}. Erstdruck, \href{https://images.p.sub.uni-goettingen.de/iiif/image/gdz:PPN319721507_1811:00000282/full/1650,/0/default.jpg}{S.~250}, \href{https://images.p.sub.uni-goettingen.de/iiif/image/gdz:PPN319721507_1811:00000283/full/1650,/0/default.jpg}{S.~251}, \href{https://images.p.sub.uni-goettingen.de/iiif/image/gdz:PPN319721507_1811:00000284/full/1650,/0/default.jpg}{S.~252}, \href{https://images.p.sub.uni-goettingen.de/iiif/image/gdz:PPN319721507_1811:00000286/full/1650,/0/default.jpg}{S.~254}, \href{https://images.p.sub.uni-goettingen.de/iiif/image/gdz:PPN319721507_1811:00000288/full/1650,/0/default.jpg}{S.~256}. \href{https://doi.org/10.3931/e-rara-73147}{Bessels Buch}.
\par\smallskip\noindent Übersetzung und Satz mit KI erstellt.
\clearpage
